\documentclass[10pt,a4paper]{amsart}
\usepackage[margin=19mm]{geometry}
\usepackage{amsmath,amssymb,mathtools}
\usepackage[scr=rsfs,cal=euler]{mathalfa}
\usepackage{xfrac}
\usepackage[unicode,hidelinks,bookmarks=false]{hyperref}
\newcommand{\ie}{i.e.\ }
\newcommand{\cf}{cf.\ }
\newcommand{\resp}{respectively\ }
\newcommand{\Subsection}{\subsection}
\providecommand{\nobreakdash}{\nobreak\hbox{-}}
\setlength{\emergencystretch}{2em}
\multlinegap0pt
\usepackage{bm}
\usepackage{bbm}
\usepackage{dsfont}
\usepackage{xspace}
\usepackage{cancel}
\usepackage{mathtools}
\usepackage{cleveref}
\usepackage[shortlabels]{enumitem}
\usepackage{amsthm}
\makeatletter
\@ifundefined{maintheorem}{\newtheorem{maintheorem}{Theorem}}{}
\@ifundefined{theorem}{\newtheorem{theorem}{Theorem}}{}
\@ifundefined{corollary}{\newtheorem{corollary}[theorem]{Corollary}}{}
\@ifundefined{lemma}{\newtheorem{lemma}[theorem]{Lemma}}{}
\@ifundefined{proposition}{\newtheorem{proposition}[theorem]{Proposition}}{}
\makeatother
\DeclareMathOperator{\PL}{PL}
\DeclareMathOperator{\ord}{ord}
\newcommand{\C}{\mathbb C}
\newcommand{\D}{\mathscr D}
\newcommand{\Lcal}{\mathscr L}
\newcommand{\Q}{\mathbb Q}
\newcommand{\R}{\mathbb R}
\newcommand{\Z}{\mathbb Z}
\title[The Multivariable Strong Monodromy Conjecture for Plane Curv]{The Multivariable Strong Monodromy Conjecture for Plane Curves}
\author{Sheng Tan}
\date{Source: arXiv:2608.26087v3 (CC BY 4.0)}
\begin{document}
\maketitle

\noindent\emph{Source and licence.} The Multivariable Strong Monodromy Conjecture for Plane Curves. Version arXiv:2608.26087v3 (CC BY 4.0), distributed under the Creative Commons Attribution 4.0 International licence, as recorded in the arXiv licence metadata for this exact version and shown on the abstract page https://arxiv.org/abs/2608.26087v3. Full rights notice: licenses/arxiv-2608.26087-CC-BY-4.0.txt.\par\smallskip

\noindent\textit{Editorial scope.} Counted unit: the Theorem whose source label is \texttt{thm:main-thm-top} in the article (source0.tex lines 2001--2006, source label \texttt{thm:main-thm-top}), delivered together with the article's own complete original proof of it (source0.tex lines 2008--2070, one \texttt{proof} environment of 63 lines). No mathematics is written, repaired or re-derived: statement and proof are the article's own text line for line. Required context retained uncounted: thm:main-intro (lines 227--232); eq:surface-zeta (lines 549--556); lem:positive-wall-normalization (lines 646--660); thm:same-wall-positivity (lines 754--769); thm:selected-component-classification (lines 832--834); thm:grouped-wall-dichotomy (lines 864--870); prop:affine-point-lifting (lines 915--924); eq:general-positive-dual-wall (lines 952--954); lem:bounded-polytope-density (lines 961--968); lem:selected-rupture-density (lines 986--997); thm:dense-wall-lifting (lines 1008--1017); thm:strict-transform-wall (lines 1033--1045); thm:direct-residue-obstruction (lines 1394--1403); prop:residue-one-injection (lines 1473--1483); thm:direct-double-residue-obstruction (lines 1543--1555); prop:multiplicity-transport (lines 1670--1680); eq:blanco-effective-punctures (lines 1808--1810); thm:blanco-residue-class (lines 1815--1838); prop:selected-numerical-exhaustion (lines 1925--1931); cor:selected-blanco-class (lines 1976--1985). Every definition, display and label of those lines is retained unchanged. Omitted from this package and not redistributed: 1--226, 233--548, 557--645, 661--753, 770--831, 835--863, 871--914, 925--951, 955--960, 969--985, 998--1007, 1018--1032, 1046--1393, 1404--1472, 1484--1542, 1556--1669, 1681--1807, 1811--1814, 1839--1924, 1932--1975, 1986--2000, 2007--2007, 2071--3188. Nothing inside a retained excerpt is omitted or altered: every retained body is delivered exactly as the article's lines are written, so no omission receipt of any kind accompanies this package. Citations: the retained excerpts carry no citation command at all, so no bibliography entry of the article is transcribed and no reference is redistributed. Reference adaptation: none was needed and none was made; every reference inside the delivered unit resolves inside it, so no unresolved reference or citation is produced. Printed slips kept literal: no printed word, symbol, hypothesis or number of the counted statement or of its proof was altered. Layout and class: the article's own class is \texttt{amsart} with packages including fontenc, inputenc, amsmath, amssymb, amsthm, mathtools, mathrsfs, libertine, newtxmath, microtype, tikz, geometry, booktabs, enumitem; the wrapper is the lane's pinned \texttt{amsart} editorial wrapper at 10pt on A4, which restates the theorem-like environments the retained text needs and adds the article's own macro definitions that the retained text needs (\texttt{\textbackslash C}, \texttt{\textbackslash D}, \texttt{\textbackslash Lcal}, \texttt{\textbackslash PL}, \texttt{\textbackslash Q}, \texttt{\textbackslash R}, \texttt{\textbackslash Z}, \texttt{\textbackslash ord}), with extra wrapper packages (bm, bbm, dsfont, xspace, cancel, mathtools, cleveref, enumitem). The delivered numbering is the unit's own. Assets: no retained span contains a figure environment, a graphics-inclusion command, a PDF-page inclusion, a table or a TikZ picture; no image file is copied into this package, the package declares no asset dependency and no image file is redistributed with it. Licence: version arXiv:2608.26087v3 of this article is distributed under the Creative Commons Attribution 4.0 International licence, as recorded in the arXiv licence metadata for this exact version and shown on the abstract page https://arxiv.org/abs/2608.26087v3. A full rights notice is delivered at licenses/arxiv-2608.26087-CC-BY-4.0.txt. Characters: every retained excerpt is pure ASCII, so the delivered engine, XeLaTeX with its default Latin Modern text font, reports no missing character.
\begin{maintheorem}[{=\,\Cref{thm:main-thm-top}}]\label{thm:main-intro}
Let $F=(f_1,\ldots,f_r)$ be a tuple of nonzero nonunit plane curve germs on a smooth complex surface germ. Then
\[
\PL\bigl(Z^{\mathrm{top}}_{F,0}\bigr)\subseteq Z(B_{F,0}).
\]
\end{maintheorem}

\begin{equation}\label{eq:surface-zeta}
 Z^{\mathrm{top}}_{F,0}(s)
 =
 \sum_{A\in J}\frac{m_A}{L_A(s)}
 +
 \sum_{\{A,B\}\subseteq J}
 \frac{m_{AB}}{L_A(s)L_B(s)}.
\end{equation}

\begin{lemma}\label{lem:positive-wall-normalization}
There is a unique equation
\begin{equation}\label{eq:normalized-wall}
 L_H(s)=1+\lambda_H\cdot s
\end{equation}
for $H$. For every $A\in S_H$ one has the exact equality of affine polynomials
\begin{equation}\label{eq:positive-wall-multiple}
 L_A=\nu_A L_H.
\end{equation}
In particular,
\[
 \lambda_H=\frac{N_A}{\nu_A}
\]
is independent of $A\in S_H$, and every proportionality coefficient $c_A:=\nu_A$ is positive.
\end{lemma}

\begin{theorem}
\label{thm:same-wall-positivity}
Suppose that $A,B\in S_H$ meet at a point of $\pi^{-1}(0)$. Then
\begin{equation}\label{eq:positive-double-coefficient}
    C_{-2}(H) = \sum_{\substack{\{A',B'\}\subseteq S_H}}\frac{m_{A'B'}}{\nu_{A'}\nu_{B'}}>0.
\end{equation}
Consequently $H$ is an actual pole of order two. Moreover, at every positive dual point
\begin{equation}\label{eq:positive-dual-wall}
    \alpha\in W_H^+ := \{\alpha\in\R_{>0}^r:\lambda_H\cdot\alpha=1\},
\end{equation}
every $A\in S_H$ satisfies 
\begin{equation} \label{eq:same-wall-zero-index}
N_A\cdot\alpha=\nu_A.
\end{equation}

\end{theorem}

\begin{theorem}\label{thm:selected-component-classification}
Assume $m_{AB}=0$ for all $A, B\in S_H$. If $A\in S_H$ satisfies $R_{A,H}\neq 0$, then $A$ is either a strict transform or an exceptional component of valency at least three. 
\end{theorem}

\begin{theorem}\label{thm:grouped-wall-dichotomy} 
Let $H$ be an irreducible component of the polar locus $\PL(Z^{\mathrm{top}}_{F,0})$. Exactly one of the following holds.
\begin{enumerate}[label=\textup{(\roman*)}]
\item Two components $A,B\in S_H$ meet at a point of $\pi^{-1}(0)$. Then $H$ is a pole of order two and the leading coefficient $C_{-2}(H)>0$. Moreover, for every $\alpha\in W_H^+$ and every $A\in S_H$, one has $N_A\cdot\alpha=\nu_A$.
\item Any two components in $S_H$ are disjoint over $\pi^{-1}(0)$. Then $H$ is a simple pole and there exists $A\in S_H$ such that $R_{A,H}\neq0$. Every such $A$ is either a strict transform or an exceptional component of valency at least three.
\end{enumerate}
\end{theorem}

\begin{proposition}\label{prop:affine-point-lifting}
Let $\alpha\in\Q_{>0}^r$. If
\begin{equation}\label{eq:cyclic-noncontainment}
 u_\alpha\notin\D_{X,0}(hu_\alpha),
\end{equation}
then
\begin{equation}\label{eq:point-in-bernstein-zero-set}
 -\alpha\in Z(B_{F,0}).
\end{equation}
\end{proposition}

\begin{equation}\label{eq:general-positive-dual-wall}
 H=V(1+\lambda\cdot s), \quad W_H=\{\alpha\in\C^r:\lambda\cdot\alpha=1\}, \quad W_H^+=W_H\cap\R_{>0}^r,
\end{equation}

\begin{lemma}\label{lem:bounded-polytope-density}
Let $H$ and $W_H^+$ be as in \eqref{eq:general-positive-dual-wall}. Let $q_1,\ldots,q_a$ be nonzero rational functions on $W_H$ defined over $\Q$, and let $\eta_1,\ldots,\eta_b$ be affine functions on $W_H$ with rational coefficients. Suppose $r\ge2$. Then there exists a bounded rational relatively open polytope $\mathcal B\subset W_H^+$ on which every $q_i$ is defined and nonzero. For each nonconstant $\eta_j$, remove from $\mathcal B$ all sets of the form
\begin{equation*}
\{\alpha\in\mathcal B:\eta_j(\alpha)=m\},\quad m\in\Z,
\end{equation*}
that meet $\mathcal B$. The rational points in the remaining set are Zariski dense in $W_H$.
If $r=1$, then $W_H$ consists of a single rational point, and the same conclusion is immediate.
\end{lemma}

\begin{lemma}\label{lem:selected-rupture-density}
Assume $m_{AB}=0$ for all $A,B\in S_H$, and let $E\in S_H$ satisfy $R_{E,H}\neq0$. There exists a Zariski-dense set
\begin{equation*}
\mathcal U_E(\Q)\subset W_H^+\cap\Q^r
\end{equation*}
such that, for every $\alpha\in\mathcal U_E(\Q)$,
\begin{enumerate}[label=\textup{(\roman*)}]
\item $\epsilon_D(\alpha)\neq0$ for every neighbour $D$ of $E$;
\item $R_{E,H}(-\alpha)\neq0$;
\item if the restriction of $\epsilon_D$ to $W_H$ is nonconstant, then $\epsilon_D(\alpha)\notin\Z$.
\end{enumerate}
\end{lemma}

\begin{theorem}\label{thm:dense-wall-lifting}
Let $H=V(1+\lambda\cdot s)$ be a rational affine hyperplane with $0\ne\lambda\in\Q_{\ge0}^r$. Suppose that
\[
 u_\alpha\notin\D_{X,0}(hu_\alpha)
\]
for a Zariski-dense set of rational points $\alpha\in W_H^+$. Then
\begin{equation}\label{eq:whole-bernstein-wall}
 H\subseteq Z(B_{F,0}).
\end{equation}
\end{theorem}

\begin{theorem}\label{thm:strict-transform-wall}
Let $C$ be an irreducible branch of the reduced divisor $\Delta=(h=0)$, and let
\[
 N_C=(\ord_C(f_1),\ldots,\ord_C(f_r)).
\]
Then
\begin{equation}\label{eq:strict-transform-wall}
 V(N_C\cdot s+1)
 \subseteq
 Z(B_{F,0}).
\end{equation}
This remains true with arbitrary positive multiplicities of $C$ in the entries of $F$.
\end{theorem}

\begin{theorem}\label{thm:direct-residue-obstruction}
Let $E$ be a component of the total transform and let $\alpha\in\Q_{>0}^r$ satisfy $N_E\cdot\alpha=\nu_E$. If
\begin{equation}\label{eq:nonzero-residue-class}
[R_{E,\alpha}]\neq0\quad\text{in}\quad H^1_{\mathrm{dR}}(E^\circ,\Lcal_{E,\alpha}),
\end{equation}
then
\begin{equation*}
u_\alpha\notin\D_{X,0}(hu_\alpha).
\end{equation*}
\end{theorem}

\begin{proposition}\label{prop:residue-one-injection}
Let $S_1$ be a finite set of intersection points $p\in E\cap A$, with $A\neq E$ an irreducible component of the reduced total transform satisfying $\epsilon_A(\alpha)=1$. Set
\begin{equation*}
E^{\mathrm{eff}}=E^\circ\cup S_1,\quad j:E^\circ\hookrightarrow E^{\mathrm{eff}}.
\end{equation*}
Then $\Lcal_{E,\alpha}$ extends across the points of $S_1$, and restriction induces an injection
\begin{equation}\label{eq:residue-one-injection}
j^*:H^1_{\mathrm{dR}}(E^{\mathrm{eff}},\Lcal_{E,\alpha})\hookrightarrow H^1_{\mathrm{dR}}(E^\circ,\Lcal_{E,\alpha}).
\end{equation}
In particular, a nonzero residue class on $E^{\mathrm{eff}}$ remains nonzero after restriction to $E^\circ$.
\end{proposition}

\begin{theorem}
\label{thm:direct-double-residue-obstruction}
Let $A$ and $B$ be any two components of the SNC total transform meeting at $p$. If
\begin{equation}\label{eq:two-wall-equalities}
 N_A\cdot\alpha=\nu_A,
 \quad
 N_B\cdot\alpha=\nu_B,
\end{equation}
then
\[
 u_\alpha\notin\D_{X,0}(hu_\alpha).
\]
\end{theorem}

\begin{proposition}\label{prop:multiplicity-transport}
The resolution formula, the grouped Laurent dichotomy, and the residue obstructions used in the proof of \Cref{thm:main-intro} (=\Cref{thm:main-thm-top}) remain valid for arbitrary nonnegative exponents $a_{ij}$. For the strict transform $C_j$ of $g_j$ one has
\begin{equation*}
N_{C_j}=(a_{1j},\ldots,a_{rj}),
\end{equation*}
and \Cref{thm:strict-transform-wall} applies with this full vector. For every positive integral scalarization $g_b=\prod_i f_i^{b_i}$, the reduced zero divisor and minimal embedded resolution are unchanged and
\begin{equation*}
\ord_A(g_b)=b\cdot N_A.
\end{equation*}
Consequently no entrywise reducedness hypothesis is required in \Cref{thm:main-intro} (=\Cref{thm:main-thm-top}).
\end{proposition}

\begin{equation}\label{eq:blanco-effective-punctures}
S_{\mathrm{eff}}(\alpha):=\{p_D:\epsilon_D(\alpha)\notin\Z\}=\{p_1,\ldots,p_k\},\quad E^{\mathrm{eff}}=E\setminus S_{\mathrm{eff}}(\alpha).
\end{equation}

\begin{theorem}\label{thm:blanco-residue-class}
Suppose that every integral adjacent residue equals $1$. Let $S_{\mathrm{eff}}(\alpha)=\{p_1,\ldots,p_k\}$ be the effective puncture set in \eqref{eq:blanco-effective-punctures}, and suppose that
\begin{equation*}
k\geq3,
 \quad
 \sum_{j=1}^k(1-\epsilon_j)=2.
\end{equation*}
Then the coefficient Poincar\'e residue defines a nonzero class
\begin{equation*}
[R_{E,\alpha}]
 \neq0
 \quad\text{in}\quad
 H^1_{\mathrm{dR}}
 \bigl(E^{\mathrm{eff}},\Lcal_{E,\alpha}\bigr).
\end{equation*}
Its restriction to the fully punctured curve is also nonzero:
\begin{equation*}
[R_{E,\alpha}]
 \neq0
 \quad\text{in}\quad
 H^1_{\mathrm{dR}}
 \bigl(E^\circ,\Lcal_{E,\alpha}\bigr).
\end{equation*}
\end{theorem}

\begin{proposition}\label{prop:selected-numerical-exhaustion}
At every point $\alpha\in\mathcal U_E(\Q)$, after filling the intersections for which $\epsilon_D(\alpha)=1$, at least three nonintegral residues remain. They satisfy
\begin{equation*}
\sum_j(1-\epsilon_j)=2.
\end{equation*}
Consequently the hypotheses of \Cref{thm:blanco-residue-class} hold.
\end{proposition}

\begin{corollary}\label{cor:selected-blanco-class}
At every point of $\mathcal U_E(\Q)$, the coefficient Poincar\'e residue is nonzero on the fully punctured curve:
\begin{equation*}
[R_{E,\alpha}]
 \neq0
 \quad\text{in}\quad
 H^1_{\mathrm{dR}}
 \bigl(E^\circ,\Lcal_{E,\alpha}\bigr).
\end{equation*}
\end{corollary}

\begin{theorem}\label{thm:main-thm-top}
For every tuple $F=(f_1,\ldots,f_r)$ of arbitrary nonzero nonunit plane curve germs on a smooth complex surface germ,
\begin{equation*}
\PL\bigl(Z^{\mathrm{top}}_{F,0}\bigr) \subseteq Z(B_{F,0}).
\end{equation*}
\end{theorem}

\begin{proof}
Let $H$ be an irreducible component of the polar locus
\begin{equation*}
\PL(Z^{\mathrm{top}}_{F,0}).
\end{equation*}
By the resolution formula \eqref{eq:surface-zeta}, $H=V(L_A)$ for at least one component $A$ of the total transform. Give it the normalization
\begin{equation*}
H=V(L_H),
 \quad
 L_H(s)=1+\lambda_H\cdot s,
\end{equation*}
of \Cref{lem:positive-wall-normalization}, and let $S_H$ be the set of all total transform components whose affine form defines $H$. We apply the two alternatives of \Cref{thm:grouped-wall-dichotomy}.

Suppose first that two components $A,B\in S_H$ meet at a point of $\pi^{-1}(0)$. At every positive rational dual point
\begin{equation*}
\alpha\in W_H^+\cap\Q^r,
 \quad
 \lambda_H\cdot\alpha=1,
\end{equation*}
\Cref{thm:same-wall-positivity} gives
\begin{equation*}
N_A\cdot\alpha=\nu_A,
 \quad
 N_B\cdot\alpha=\nu_B.
\end{equation*}
The double-residue obstruction \Cref{thm:direct-double-residue-obstruction} therefore yields
\begin{equation*}
u_\alpha\notin\D_{X,0}(hu_\alpha).
\end{equation*}
The rational points of $W_H^+$ are Zariski dense in $W_H$ when $r\geq2$ by \Cref{lem:bounded-polytope-density}; for $r=1$ the hyperplane is the corresponding single rational point. Hence \Cref{thm:dense-wall-lifting} gives
\begin{equation*}
H\subseteq Z(B_{F,0}).
\end{equation*}

Suppose next that any two components in $S_H$ are disjoint over $\pi^{-1}(0)$. By \Cref{thm:grouped-wall-dichotomy}, there is a component $A\in S_H$ with $R_{A,H}\neq0$. By \Cref{thm:selected-component-classification}, $A$ is either a strict transform or a rupture component. In the strict-transform case, $H=V(N_A\cdot s+1)$, and the desired inclusion follows directly from \Cref{thm:strict-transform-wall}.

It remains to treat the case $A=E$ with $E$ rupture. Let
\begin{equation*}
\mathcal U_E(\Q)\subset W_H^+\cap\Q^r
\end{equation*}
be the Zariski-dense set from \Cref{lem:selected-rupture-density}. For every $\alpha\in\mathcal U_E(\Q)$, \Cref{prop:selected-numerical-exhaustion}, \Cref{thm:blanco-residue-class}, and \Cref{prop:residue-one-injection} give
\begin{equation*}
[R_{E,\alpha}]\neq0
 \quad\text{in}\quad
 H^1_{\mathrm{dR}}(E^\circ,\Lcal_{E,\alpha});
\end{equation*}
equivalently, this is \Cref{cor:selected-blanco-class}. The residue obstruction \Cref{thm:direct-residue-obstruction} now gives
\begin{equation*}
u_\alpha\notin\D_{X,0}(hu_\alpha)
 \quad(\alpha\in\mathcal U_E(\Q)).
\end{equation*}
\Cref{prop:affine-point-lifting} places every $-\alpha$ in $Z(B_{F,0})$, and \Cref{thm:dense-wall-lifting} yields
\begin{equation*}
H\subseteq Z(B_{F,0}).
\end{equation*}

We have therefore shown that every irreducible component $H$ of $\PL(Z^{\mathrm{top}}_{F,0})$ is contained in $Z(B_{F,0})$. Taking the union over all irreducible components gives
\begin{equation*}
\PL\bigl(Z^{\mathrm{top}}_{F,0}\bigr)
  \subseteq Z(B_{F,0}).
\end{equation*}
\Cref{prop:multiplicity-transport} shows that the same argument covers powers, repeated branches, common factors, and dependent tuple entries. The point case $r=1$ was included in both density steps above.
\end{proof}

\end{document}
